% GENERATED FILE. Edit canonical lessons, then run npm run generate:books.
% canonical-source-hash: fnv1a64:7c3a67fdeeb1de63
% canonical-lessons: ES-C635-pacto, ES-C635-cuestionar, ES-C635-avalar, ES-C635-coherencia, ES-C635-dispuesto

\chapter{Pactos y dudas}
\label{ch:es-pactos-y-dudas}
\hlchaptermodality{635}

\begin{chapteropening}
\textbf{By the end of this chapter:} \emph{I can make a deal, vouch for someone, question what I doubt and say whether ideas are consistent.}

Complete the last lesson of chapter 635: I can make a deal, vouch for someone, question what I doubt and say whether ideas are consistent.
\end{chapteropening}

\section[el pacto]{el pacto — the pact, the deal}

\label{lesson:ES-C635-pacto}



\begin{quote}
Before the new one: say the Spanish for to agree on, then the Spanish for to persuade.
\end{quote}

\subsection*{You'll want to know: el pacto}
\textbf{el pacto} — \textquotedblleft{}the pact, the deal\textquotedblright{}.

\begin{quote}
Hicimos un pacto: yo cocino y tú limpias.
\end{quote}

\begin{quote}
\emph{We made a deal: I cook and you clean.}
\end{quote}

\emph{Pactar} — to agree a deal.

\begin{cousinweb}
\emph{pacto} is a learned borrowing from Latin \emph{pactum} 'an agreement' (EtymDB), from \emph{pacisci} 'to make a bargain'.
\end{cousinweb}

\begin{grammarlens}[title={What you've built}]
One more word for telling, explaining and describing at B1.
\end{grammarlens}

\subsection*{Your turn}
\begin{itemize}
\raggedright
  \item \emph{Say it:} \emph{el pacto}
  \item \emph{Say it:} \emph{el pacto}, once more
  \item \emph{Say it:} a sentence of your own with \emph{el pacto}
  \item \emph{Recall:} say the Spanish for to agree on, then the Spanish for to persuade, then say \emph{el pacto} again
\end{itemize}

\begin{tcolorbox}[breakable,colback=teal!4,colframe=teal!35!black,title={Before you move on}]
What does el pacto mean? (\textquotedblleft{}The pact, the deal\textquotedblright{}.) Say it once more.
\end{tcolorbox}
\section[cuestionar]{cuestionar — to question, to call into doubt}

\label{lesson:ES-C635-cuestionar}



\begin{quote}
Before the new one: say the Spanish for to persuade, then the Spanish for the pact.
\end{quote}

\subsection*{You'll want to know: cuestionar}
\textbf{cuestionar} — \textquotedblleft{}to question, to call into doubt\textquotedblright{}.

\begin{quote}
Nadie cuestiona su palabra.
\end{quote}

\begin{quote}
\emph{Nobody questions his word.}
\end{quote}

It means to doubt something, not to ask a question: that is \emph{preguntar}.

\begin{cousinweb}
\emph{cuestionar} is built on \emph{cuestión}, a learned borrowing from Latin \emph{quaestio} 'inquiry, question', from \emph{quaerere} 'to seek, to ask' (EtymDB).
\end{cousinweb}

\begin{grammarlens}[title={What you've built}]
One more word for telling, explaining and describing at B1.
\end{grammarlens}

\subsection*{Your turn}
\begin{itemize}
\raggedright
  \item \emph{Say it:} \emph{cuestionar}
  \item \emph{Say it:} \emph{cuestionar}, once more
  \item \emph{Say it:} a sentence of your own with \emph{cuestionar}
  \item \emph{Recall:} say the Spanish for to persuade, then the Spanish for the pact, then say \emph{cuestionar} again
\end{itemize}

\begin{tcolorbox}[breakable,colback=teal!4,colframe=teal!35!black,title={Before you move on}]
What does cuestionar mean? (\textquotedblleft{}To question, to call into doubt\textquotedblright{}.) Say it once more.
\end{tcolorbox}
\section[avalar]{avalar — to vouch for, to guarantee; to back}

\label{lesson:ES-C635-avalar}



\begin{quote}
Before the new one: say the Spanish for the pact, then the Spanish for to question.
\end{quote}

\subsection*{You'll want to know: avalar}
\textbf{avalar} — \textquotedblleft{}to vouch for, to guarantee; to back\textquotedblright{}.

\begin{quote}
Mis padres me avalan para el piso.
\end{quote}

\begin{quote}
\emph{My parents are my guarantors for the flat.}
\end{quote}

\emph{El aval} — the guarantee; \emph{el avalista} — the guarantor.

\begin{cousinweb}
\emph{avalar} is built on \emph{aval} 'guarantee', which was borrowed from French \emph{aval} (EtymDB); where the French word comes from is not settled.
\end{cousinweb}

\begin{grammarlens}[title={What you've built}]
One more word for telling, explaining and describing at B1.
\end{grammarlens}

\subsection*{Your turn}
\begin{itemize}
\raggedright
  \item \emph{Say it:} \emph{avalar}
  \item \emph{Say it:} \emph{avalar}, once more
  \item \emph{Say it:} a sentence of your own with \emph{avalar}
  \item \emph{Recall:} say the Spanish for the pact, then the Spanish for to question, then say \emph{avalar} again
\end{itemize}

\begin{tcolorbox}[breakable,colback=teal!4,colframe=teal!35!black,title={Before you move on}]
What does avalar mean? (\textquotedblleft{}To vouch for, to guarantee; to back\textquotedblright{}.) Say it once more.
\end{tcolorbox}
\section[la coherencia]{la coherencia — the consistency, the coherence}

\label{lesson:ES-C635-coherencia}



\begin{quote}
Before the new one: say the Spanish for to question, then the Spanish for to vouch for.
\end{quote}

\subsection*{You'll want to know: la coherencia}
\textbf{la coherencia} — \textquotedblleft{}the consistency, the coherence\textquotedblright{}.

\begin{quote}
Sus ideas tienen coherencia.
\end{quote}

\begin{quote}
\emph{Her ideas are consistent.}
\end{quote}

\emph{Coherente} — consistent.

\begin{cousinweb}
\emph{coherencia} is a learned borrowing from Latin \emph{cohaerentia} (EtymDB), built on \emph{cohaerere} 'to stick together', from \emph{haerere} 'to stick'.
\end{cousinweb}

\begin{grammarlens}[title={What you've built}]
One more word for telling, explaining and describing at B1.
\end{grammarlens}

\subsection*{Your turn}
\begin{itemize}
\raggedright
  \item \emph{Say it:} \emph{la coherencia}
  \item \emph{Say it:} \emph{la coherencia}, once more
  \item \emph{Say it:} a sentence of your own with \emph{la coherencia}
  \item \emph{Recall:} say the Spanish for to question, then the Spanish for to vouch for, then say \emph{la coherencia} again
\end{itemize}

\begin{tcolorbox}[breakable,colback=teal!4,colframe=teal!35!black,title={Before you move on}]
What does la coherencia mean? (\textquotedblleft{}The consistency, the coherence\textquotedblright{}.) Say it once more.
\end{tcolorbox}
\section[dispuesto]{dispuesto — willing, ready}

\label{lesson:ES-C635-dispuesto}



\begin{quote}
Before the new one: say the Spanish for to vouch for, then the Spanish for the consistency.
\end{quote}

\subsection*{You'll want to know: dispuesto}
\textbf{dispuesto} — \textquotedblleft{}willing, ready\textquotedblright{}.

\begin{quote}
Estoy dispuesto a ayudar.
\end{quote}

\begin{quote}
\emph{I am willing to help.}
\end{quote}

\emph{Estar dispuesto a} + infinitive.

\begin{cousinweb}
\emph{dispuesto} is the participle of \emph{disponer}, a learned borrowing from Latin \emph{disponere} 'to set in order', built from \emph{dis} and \emph{ponere} 'to put' (EtymDB).
\end{cousinweb}

\begin{grammarlens}[title={What you've built}]
One more word for telling, explaining and describing at B1.
\end{grammarlens}

\subsection*{Your turn}
\begin{itemize}
\raggedright
  \item \emph{Say it:} \emph{dispuesto}
  \item \emph{Say it:} \emph{dispuesto}, once more
  \item \emph{Say it:} a sentence of your own with \emph{dispuesto}
  \item \emph{Recall:} say the Spanish for to vouch for, then the Spanish for the consistency, then say \emph{dispuesto} again
\end{itemize}

\begin{tcolorbox}[breakable,colback=teal!4,colframe=teal!35!black,title={Before you move on}]
What does dispuesto mean? (\textquotedblleft{}Willing, ready\textquotedblright{}.) Say it once more.
\end{tcolorbox}
